% THIS DOCUMENT IS FOLLOWS THE VOLERE TEMPLATE BY Suzanne Robertson and James Robertson
% ONLY THE SECTION HEADINGS ARE PROVIDED
%
% Initial draft from https://github.com/Dieblich/volere
%
% Risks are removed because they are covered by the Hazard Analysis
\documentclass[12pt]{article}

\usepackage{booktabs}
\usepackage{tabularx}
\usepackage{enumitem}
\usepackage{float}
\usepackage{hyperref}
\hypersetup{
    bookmarks=true,         % show bookmarks bar?
      colorlinks=true,      % false: boxed links; true: colored links
    linkcolor=red,          % color of internal links (change box color with linkbordercolor)
    citecolor=green,        % color of links to bibliography
    filecolor=magenta,      % color of file links
    urlcolor=cyan           % color of external links
}

\newcommand{\lips}{\textit{Insert your content here.}}

\input{../Comments.text}
\input{../Common.text}

\begin{document}

\title{Software Requirements Specification for \progname: subtitle describing software} 
\author{\authname}
\date{\today}
	
\maketitle

~\newpage

\pagenumbering{roman}

\tableofcontents

~\newpage

\section*{Revision History}

\begin{tabularx}{\textwidth}{p{3cm}p{2cm}X}
\toprule {\textbf{Date}} & {\textbf{Version}} & {\textbf{Notes}}\\
\midrule
Date 1 & 1.0 & Notes\\
Date 2 & 1.1 & Notes\\
\bottomrule
\end{tabularx}

~\\

~\newpage
\section{Purpose of the Project}
\subsection{User Business}
\lips
\subsection{Goals of the Project}
\lips
\section{Stakeholders}
\subsection{Client}
\lips
\subsection{Customer}
\lips
\subsection{Other Stakeholders}
\lips
\subsection{Hands-On Users of the Project}
\lips
\subsection{Personas}
\lips
\subsection{Priorities Assigned to Users}
\lips
\subsection{User Participation}
\lips
\subsection{Maintenance Users and Service Technicians}
\lips

\section{Mandated Constraints}
\subsection{Solution Constraints}
\lips
\subsection{Implementation Environment of the Current System}
\lips
\subsection{Partner or Collaborative Applications}
\lips
\subsection{Off-the-Shelf Software}
\lips
\subsection{Anticipated Workplace Environment}
\lips
\subsection{Schedule Constraints}
\lips
\subsection{Budget Constraints}
\lips
\subsection{Enterprise Constraints}
\lips

\section{Naming Conventions and Terminology}
\subsection{Glossary of All Terms, Including Acronyms, Used by Stakeholders
involved in the Project}
\lips

\section{Relevant Facts And Assumptions}
\subsection{Relevant Facts}
\lips
\subsection{Business Rules}
\lips
\subsection{Assumptions}
\lips

\section{The Scope of the Work}
\subsection{The Current Situation}
\lips
\subsection{The Context of the Work}
\lips
\subsection{Work Partitioning}
\lips
\subsection{Specifying a Business Use Case (BUC)}
\lips

\section{Business Data Model and Data Dictionary}
\subsection{Business Data Model}
\lips
\subsection{Data Dictionary}
\lips

\section{The Scope of the Product}
\subsection{Product Boundary}
\lips
\subsection{Product Use Case Table}
\lips
\subsection{Individual Product Use Cases (PUC's)}
\lips

\section{Functional Requirements}
\subsection{Functional Requirements}
\lips

\section{Look and Feel Requirements}
\subsection{Appearance Requirements}
\lips
\subsection{Style Requirements}
\lips

\section{Usability and Humanity Requirements}
\subsection{Ease of Use Requirements}
\lips
\subsection{Personalization and Internationalization Requirements}
\lips
\subsection{Learning Requirements}
\lips
\subsection{Understandability and Politeness Requirements}
\lips
\subsection{Accessibility Requirements}
\lips

\section{Performance Requirements}
\subsection{Speed and Latency Requirements}
\lips
\subsection{Safety-Critical Requirements}
\lips
\subsection{Precision or Accuracy Requirements}
\lips
\subsection{Robustness or Fault-Tolerance Requirements}
\lips
\subsection{Capacity Requirements}
\lips
\subsection{Scalability or Extensibility Requirements}
\lips
\subsection{Longevity Requirements}
\lips

\section{Operational and Environmental Requirements}
\subsection{Expected Physical Environment}
\lips
\subsection{Wider Environment Requirements}
\lips
\subsection{Requirements for Interfacing with Adjacent Systems}
\lips
\subsection{Productization Requirements}
\lips
\subsection{Release Requirements}
\lips

\section{Maintainability and Support Requirements}
\subsection{Maintenance Requirements}
\lips
\subsection{Supportability Requirements}
\lips
\subsection{Adaptability Requirements}
\lips

\section{Security Requirements}
\subsection{Access Requirements}
\lips
\subsection{Integrity Requirements}
\lips
\subsection{Privacy Requirements}
\lips
\subsection{Audit Requirements}
\lips
\subsection{Immunity Requirements}
\lips

\section{Cultural Requirements}
\subsection{Cultural Requirements}
\lips

\section{Compliance Requirements}
\subsection{Legal Requirements}
\lips
\subsection{Standards Compliance Requirements}
\lips

\section{Open Issues}
\lips

\section{Off-the-Shelf Solutions}
\subsection{Ready-Made Products}
\lips
\subsection{Reusable Components}
\lips
\subsection{Products That Can Be Copied}
\lips

\section{New Problems}
\subsection{Effects on the Current Environment}
\lips
\subsection{Effects on the Installed Systems}
\lips
\subsection{Potential User Problems}
\lips
\subsection{Limitations in the Anticipated Implementation Environment That May
Inhibit the New Product}
\lips
\subsection{Follow-Up Problems}
\lips

\section{Tasks}
\subsection{Project Planning}
\lips
\subsection{Planning of the Development Phases}
\lips

\section{Migration to the New Product}
\subsection{Requirements for Migration to the New Product}
\lips
\subsection{Data That Has to be Modified or Translated for the New System}
\lips

\section{Costs}
\lips
\section{User Documentation and Training}
\subsection{User Documentation Requirements}
\lips
\subsection{Training Requirements}
\lips

\section{Waiting Room}
\lips

\section{Ideas for Solution}
\lips

\newpage{}
\section{Appendix --- Requirement Sources}
\label{app:sources}

This appendix is the requirement sources register. Each source has a fixed
identifier (SRC-001 to SRC-007) that is never renumbered. A requirement names
the sources it comes from in its rationale or source field. A statement from
the specialist interview is cited by its source identifier, participant code
and transcript timestamp, for example ``SRC-005 (P1, 28:37)''; the statement
identifiers in Section~\ref{sec:src-interview} (for example SRC-005.01) can
be cited instead. The interview is expert-reported evidence only. The three
participants are de-identified domain specialists (P1, P2 and P3). No child
took part in the interview or in any other source listed here.

Table~\ref{tab:src-docs} lists the project documents and meeting sources, and
Table~\ref{tab:src-other} lists the interview and peer-review sources. The
supported requirement identifiers are not listed in these tables yet. They
are added as back-links as each requirement is written, so that every
dependency appears at both ends: the requirement names its source, and the
source lists the requirement. Each source and statement has a \LaTeX{} label
(for example \texttt{src:005.01}), so a requirement can link to it and the
link works in the PDF.

The sources represent users in different ways. Children are the main
beneficiaries, but by the supervisor's direction (SRC-004) no requirement is
gathered from children and no information about them is collected. Their
needs enter the requirements only through adults who know them. The domain
specialists (SRC-005) are the only outside voice so far. Their advice changed
the project: progression is chosen by an adult instead of adjusted
automatically, finger input is the starting point, and camera tracing became
a stretch goal. Parents, educators and occupational therapists have not been
consulted yet. The specialists advised doing so (SRC-005.03, SRC-005.06), and
until then the requirements that depend on their views rest on the expert
statements below and are open to change.

\begin{table}[H]
\centering
\caption{Requirement sources: project documents and meetings}
\label{tab:src-docs}
\footnotesize
\begin{tabularx}{\textwidth}{@{}
  >{\raggedright\arraybackslash}p{1.3cm}
  >{\raggedright\arraybackslash}p{1.9cm}
  >{\raggedright\arraybackslash}p{1.4cm}
  >{\raggedright\arraybackslash}X
  >{\raggedright\arraybackslash}X
  >{\raggedright\arraybackslash}X@{}}
\toprule
\textbf{ID} & \textbf{Type} & \textbf{Date} & \textbf{Reference} &
\textbf{Summary} & \textbf{Limitations} \\
\midrule
\phantomsection\label{src:001}SRC-001 & Project document & -- &
\texttt{docs/\allowbreak Problem\allowbreak Statement\allowbreak And\allowbreak
  Goals/\allowbreak Problem\allowbreak Statement.tex} &
Problem, stakeholders, goals and stretch goals. &
Written by the team; some assumptions changed after SRC-005. \\
\phantomsection\label{src:002}SRC-002 & Project document & -- &
\texttt{docs/\allowbreak Development\allowbreak Plan/\allowbreak
  Development\allowbreak Plan.tex} &
Team process, tools and schedule. &
Written by the team; plans, not user needs. \\
\phantomsection\label{src:003}SRC-003 & TA meeting & Sep 23, 2026
  & Issue \href{https://github.com/AHMEDELZARIA/squigglr/issues/10}{\#10} &
Website instead of an installed app; teachers and students as primary
stakeholders; specific, measurable metrics. &
Short minutes; not a user source. \\
\phantomsection\label{src:004}SRC-004 & Supervisor meeting & Sep 23, 2026
  & Issue \href{https://github.com/AHMEDELZARIA/squigglr/issues/21}{\#21} &
Collect no information about children; do not gather requirements by asking
children; ethics approval for studies; start with a web app; traceability. &
Advice from the client, not observed user data. \\
\bottomrule
\end{tabularx}
\end{table}

\begin{table}[H]
\centering
\caption{Requirement sources: interview and reviews}
\label{tab:src-other}
\footnotesize
\begin{tabularx}{\textwidth}{@{}
  >{\raggedright\arraybackslash}p{1.3cm}
  >{\raggedright\arraybackslash}p{1.9cm}
  >{\raggedright\arraybackslash}p{1.4cm}
  >{\raggedright\arraybackslash}X
  >{\raggedright\arraybackslash}X
  >{\raggedright\arraybackslash}X@{}}
\toprule
\textbf{ID} & \textbf{Type} & \textbf{Date} & \textbf{Reference} &
\textbf{Summary} & \textbf{Limitations} \\
\midrule
\phantomsection\label{src:005}SRC-005 & Specialist interview & Oct 1, 2026 &
Appendix~\ref{app:interview} &
Three de-identified specialists on progression, input method, child-friendly
design, feedback, accessibility and ethics. &
Three experts, one session; expert opinion, not child observation; the
experts advised confirming details with occupational therapists and
educators. \\
\phantomsection\label{src:006}SRC-006 & TA meeting & Oct 7, 2026
  & Issue \href{https://github.com/AHMEDELZARIA/squigglr/issues/78}{\#78} &
State machine for the formal specification; define jargon; named TBD
constants; specialists as de-identified secondary stakeholders. &
Course guidance, not user needs. \\
\phantomsection\label{src:007}SRC-007 & Peer review & --
  & Issues \href{https://github.com/AHMEDELZARIA/squigglr/issues/69}{\#69}--%
  \href{https://github.com/AHMEDELZARIA/squigglr/issues/75}{\#75} &
Team 13 review of SRC-001 and SRC-002. &
Reviewers' opinions; not all suggestions are adopted. \\
\bottomrule
\end{tabularx}
\end{table}

\subsection{Key Interview Statements}
\label{sec:src-interview}

This subsection lists the interview statements most relevant to the
requirements. Each is a short paraphrase, not a quote; the exact words are in
Appendix~\ref{app:interview} at the given timestamp.

\begin{description}[leftmargin=1.5em, style=nextline]
\item[\phantomsection%
  \label{src:005.01}SRC-005.01 (P1, 28:37; P2, 29:18; 30:20--31:02)]
Deciding automatically when a child is ready for the next level would need
assessment and is too complex for the time frame. An adult who knows the
child (a parent or educator) should choose the level. The team agreed to
leave progression to the supervising adult.
\item[\phantomsection%
  \label{src:005.02}SRC-005.02 (P2, 29:18)]
In a past project, a teacher personalized each child's activities through a
separate adult interface that also showed progress.
\item[\phantomsection%
  \label{src:005.03}SRC-005.03 (P2, 31:42)]
Which progress measures help an adult should be decided with occupational
therapists or educational specialists, using milestones from the literature
if they exist.
\item[\phantomsection%
  \label{src:005.04}SRC-005.04 (P3, 24:34; P3, 25:36; P1, 26:46)]
Holding a pen affects how children draw and the usability of the app.
Finger input is the easiest start; a larger stylus and then a smaller one
could be later difficulty levels if fine motor practice is a goal.
\item[\phantomsection%
  \label{src:005.05}SRC-005.05 (P1, 34:13)]
The app should be child-friendly: a child who cannot read yet should be able
to enter a lesson, practise, restart and move back and forth using images
and icons rather than text, without constant adult help. Adults handle setup.
\item[\phantomsection%
  \label{src:005.06}SRC-005.06 (P1, 36:26)]
Positive feedback wording and colours should be designed with feedback from
educators and occupational therapists, built up bit by bit.
\item[\phantomsection%
  \label{src:005.07}SRC-005.07 (P2, 37:39)]
Audio feedback is a good idea; offering more than one modality is
encouraged in inclusive design, with the user able to turn the audio off.
Design it carefully with therapists' feedback, and consider adults with ADHD
or autism as reviewers when children cannot take part.
\item[\phantomsection%
  \label{src:005.08}SRC-005.08 (P2, 39:26)]
Start with the Web Content Accessibility Guidelines (WCAG), then apply more
specialized design heuristics for autism from the literature.
\item[\phantomsection%
  \label{src:005.09}SRC-005.09 (P2, 19:05; P1, 20:27; P2, 20:47)]
From a universal design view, the skills apply to all children, including
those with motor or sight challenges; designing for disability serves many
groups.
\item[\phantomsection%
  \label{src:005.10}SRC-005.10 (P1, 13:42; P1, 15:12; P1, 16:13; P2, 16:45)]
Camera finger tracing in the air practises gross rather than fine motor
skills. It is useful but should come after the core touchscreen tracing
when time is limited.
\item[\phantomsection%
  \label{src:005.11}SRC-005.11 (P2, 22:56; P1, 24:03; P1, 33:00)]
Testing with children needs ethics approval (about three to four months
where P2 works). Meeting experts for feedback without involving children
does not.
\item[\phantomsection%
  \label{src:005.12}SRC-005.12 (41:51--42:48)]
All three participants agreed to the use of the transcript for eliciting
requirements.
\end{description}

\newpage{}
\input{InterviewTranscript.text}

\newpage{}
\section{Appendix --- Generative AI Use Disclosure}

\wss{ChatGPT (version 5) was used to brainstorm ideas for the template for this disclose section.}

\subsection{AI Tools Used}

\wss{Give the name of the AI tool(s) used and the version.}

\subsection{How and Where Used}

\wss{Explain what the tool(s) were used for.  Examples include: brainstorming,
explaining a technical concept, reviewing the document, generating or modifying
text, generating or modifying diagrams, identifying errors or omissions or
inconsistencies.}

\wss{For each of the uses, identify which parts of the document were affected.}

\wss{You do not need to report every interaction, but you should disclose the 
uses that materially contributed to your work}

\subsection{Verification of AI-Generated Content}

\wss{\begin{itemize}
    \item Describe how the team checked AI-generated or AI-assisted material
    before including it in the document.
    \item In particular, verify factual claims, technical assertions,
    requirements, calculations, references, and other information for which an
    AI system may produce incorrect or fabricated results.
    \item \textbf{The team is responsible for the accuracy, completeness,
consistency, and appropriateness of the submitted document, regardless of
whether material was generated by a human or an AI system.}
\end{itemize}}

\newpage{}

\section{Appendix --- Reflection}

\input{../Reflection.text}

\input{../SRS_Reflection.text}

\end{document}
